% ---
% title: Pairwise but not mutually independent
% tags: independence
% ---

\begin{problem}
	Two fair coins are tossed. Let $A$ = ``the first coin shows heads'', $B$ =
	``the second coin shows heads'' and $C$ = ``the two coins show the same
	side''. (a) Are $A$ and $B$ independent? $A$ and $C$? $B$ and $C$? (b) Are
	$A$, $B$, $C$ (mutually) independent? (c) Explain in words why knowing two
	of the events tells you everything about the third.
\end{problem}

\begin{slides}
	\small
	\exstatement
	\pause
	\begin{columns}[T]
		\begin{column}{.6\textwidth}
			\sol $\Om=\{HH,HT,TH,TT\}$, four equally likely outcomes,
			each with probability $\frac14$.
			\pause

			\smallskip
			$A=\{HH,HT\}$, \ $B=\{HH,TH\}$, \ $C=\{HH,TT\}$, so
			$P(A)=P(B)=P(C)=\frac24=\frac12$.
		\end{column}
		\begin{column}{.36\textwidth}
			\uncover<3->{
				\centering
				\begin{tabular}{c|ccc}
					     & $A$ & $B$ & $C$ \\ \hline
					$HH$ & $\bullet$ & $\bullet$ & $\bullet$ \\
					$HT$ & $\bullet$ &            &            \\
					$TH$ &            & $\bullet$ &            \\
					$TT$ &            &            & $\bullet$
				\end{tabular}
			}
		\end{column}
	\end{columns}
	\pause
	\smallskip
	(a) Each pair has only $HH$ in common:
	$AB=AC=BC=\{HH\}$.
	\pause\\[1mm]
	$P(AB)=\frac14=\frac12\cdot\frac12=P(A)P(B)$; in the same way
	$P(AC)=P(A)P(C)$ and $P(BC)=P(B)P(C)$.
	\pause\\[1mm]
	\alert{Every pair is independent} (``pairwise independent'').
\end{slides}

\begin{slides}
	\small
	Three events $E,F,G$ are \alert{independent} if \emph{all four}
	equations hold:
	\[ P(EF)=P(E)P(F),\quad P(EG)=P(E)P(G),\quad P(FG)=P(F)P(G), \]
	\[ P(EFG)=P(E)P(F)P(G). \]
	\pause
	(b) The first three hold by (a). The fourth: $ABC=\{HH\}$, so
	\[ P(ABC)=\tfrac14, \qquad\text{but}\qquad
		P(A)P(B)P(C)=\tfrac12\cdot\tfrac12\cdot\tfrac12=\tfrac18. \]
	\pause
	$\tfrac14\ne\tfrac18$: $A$, $B$, $C$ are \alert{not independent}.
	\pause

	\medskip
	(c) If $A$ and $B$ both occur, the outcome is $HH$, so $C$ occurs:
	\[ P(C\mid AB)=1\ne\tfrac12=P(C). \]
	\pause
	One coin alone says nothing about ``same side''; two of the events
	together decide the third.
	\pause

	\medskip
	\alert{Moral:} pairwise independence is not enough; check all four
	equations.
\end{slides}

\begin{solution}
	\idea List the four equally likely outcomes, write each event as a set of
	outcomes, and compare the probability of each intersection with the
	product of the probabilities.

	\step{Sample space and events.} The coins are fair and tossed
	independently, so the four outcomes
	\[ \Om=\{HH,HT,TH,TT\} \]
	are equally likely, each with probability $\frac14$ (the first letter is
	the first coin). The events as sets:
	\[ A=\{HH,HT\},\qquad B=\{HH,TH\},\qquad C=\{HH,TT\}. \]
	Each contains two of the four outcomes, so
	$P(A)=P(B)=P(C)=\frac24=\frac12$. In the table below, a dot means that
	the event of the column occurs at the outcome of the row.

	\begin{center}
		\begin{tabular}{c|ccc}
			     & $A$ & $B$ & $C$ \\ \hline
			$HH$ & $\bullet$ & $\bullet$ & $\bullet$ \\
			$HT$ & $\bullet$ &            &            \\
			$TH$ &            & $\bullet$ &            \\
			$TT$ &            &            & $\bullet$
		\end{tabular}
	\end{center}

	\step{(a)} Two events $E$ and $F$ are independent when
	$P(EF)=P(E)P(F)$. In the table, the only outcome in two columns at once is
	$HH$. So every pair of events has the same intersection:
	\[ AB=AC=BC=\{HH\}. \]
	Each of these has probability $\frac14$, and each product of two
	probabilities is $\frac12\cdot\frac12=\frac14$:
	\[ P(AB)=\tfrac14=P(A)P(B),\quad P(AC)=\tfrac14=P(A)P(C),\quad
		P(BC)=\tfrac14=P(B)P(C). \]
	All three pairs are independent. For $A$ and $B$ this is expected (two
	separate coins). For $A$ and $C$ it is less obvious: if the first coin
	shows heads, ``same side'' means the second coin shows heads too, which
	still has probability $\frac12$; so $P(C\mid A)=\frac12=P(C)$.

	\step{(b)} Three events $E$, $F$, $G$ are independent only if \emph{all
	four} equations hold:
	\[ P(EF)=P(E)P(F),\quad P(EG)=P(E)P(G),\quad P(FG)=P(F)P(G),\quad
		P(EFG)=P(E)P(F)P(G). \]
	The first three hold by (a). For the fourth, all three events occur only
	at $HH$, so $ABC=\{HH\}$ and
	\[ P(ABC)=\frac14,\qquad\text{while}\qquad
		P(A)P(B)P(C)=\frac12\cdot\frac12\cdot\frac12=\frac18. \]
	Since $\frac14\ne\frac18$, the fourth equation fails, and $A$, $B$, $C$
	are not independent.

	\step{(c)} Suppose we know that $A$ and $B$ both occurred. Then both
	coins show heads, so they show the same side, and $C$ certainly occurred:
	\[ P(C\mid AB)=\frac{P(ABC)}{P(AB)}=\frac{1/4}{1/4}=1, \]
	while $P(C)=\frac12$. In the same way, any two of the events decide the
	third: if $A$ occurs and $C$ occurs, the second coin must be heads, so
	$B$ occurs; if $A$ occurs and $C$ does not, $B$ does not occur either.
	Knowing only one coin tells nothing about ``same side'' (that is
	the pairwise independence of (a)), but knowing both coins tells
	everything.
	\begin{answer}
		(a) yes, all three pairs are independent:
		$P(AB)=P(AC)=P(BC)=\frac14=\frac12\cdot\frac12$ \quad
		(b) no: $P(ABC)=\frac14\ne\frac18=P(A)P(B)P(C)$ \quad
		(c) $P(C\mid AB)=1\ne\frac12=P(C)$: two of the events determine the
		third
	\end{answer}
	\begin{remark}
		\textbf{Common mistake:} checking only the three pairs and concluding
		that $A$, $B$, $C$ are independent. Pairwise independence does not
		imply independence of all three; the equation for the triple
		intersection must be checked separately. (The converse fails too: the
		triple equation alone does not imply the pairwise ones.)
	\end{remark}
\end{solution}
